% GENERATED FILE. Edit canonical lessons, then run npm run generate:books.
% canonical-source-hash: fnv1a64:353ecd016d3928a8
% canonical-lessons: KA-C461-atithigrha, KA-C461-hastel, KA-C461-thevani, KA-C461-lakar, KA-C461-marupavati, KA-R461-review-455-461

\chapter{Guest Houses and Deposits}
\label{ch:ka-guest-houses-and-deposits}
\hlchaptermodality{461}

\begin{chapteropening}
\textbf{By the end of this chapter:} \emph{I can say a guest house, a hostel, a security deposit, a locker and a refund.}

Complete the last lesson of chapter 461: I can say a guest house, a hostel, a security deposit, a locker and a refund.
\end{chapteropening}

\section[\texorpdfstring{\kn{ಅತಿಥಿಗೃಹ} (atithigṛha)}{ಅತಿಥಿಗೃಹ (atithigṛha)}]{\kn{ಅತಿಥಿಗೃಹ} (atithigṛha) — a guest house}

\label{lesson:KA-C461-atithigrha}



\begin{quote}
Before the new one: say the Kannada for a receptionist, then the Kannada for a counter.
\end{quote}

\subsection*{You'll want to know: \kn{ಅತಿಥಿಗೃಹ}}
\textbf{\kn{ಅತಿಥಿಗೃಹ}} — \emph{atithigṛha} — \textquotedblleft{}a guest house\textquotedblright{}.

\begin{quote}
\kn{ಸರ್ಕಾರಿ} \kn{ಅತಿಥಿಗೃಹದಲ್ಲಿ} \kn{ಕೊಠಡಿಗಳು} \kn{ಅಗ್ಗ}.
\end{quote}

\begin{quote}
\emph{sarkāri atithigṛhadalli koṭhaḍigaḷu agga.}
\end{quote}

\begin{quote}
Rooms in the government guest house are cheap.
\end{quote}

\kn{ಅತಿಥಿ}, a guest, plus \kn{ಗೃಹ}, a house.

\begin{cousinweb}
\textbf{\kn{ಅತಿಥಿಗೃಹ}} (\emph{atithigṛha}) is built from two Sanskrit words: \emph{atithi}, \textquotedblleft{}a guest\textquotedblright{}, and \emph{gṛha}, \textquotedblleft{}a house\textquotedblright{} (Monier-Williams).
\end{cousinweb}

\begin{grammarlens}[title={What you've built}]
One more word for getting around.
\end{grammarlens}

\subsection*{Your turn}
\begin{itemize}
\raggedright
  \item \emph{Say it:} \emph{atithigṛha}
  \item \emph{Say it:} \emph{atithigṛha}, once more
  \item \emph{Say it:} \emph{sarkāri atithigṛhadalli koṭhaḍigaḷu agga.} — Rooms in the government guest house are cheap.
  \item \emph{Recall:} say the Kannada for a receptionist, then the Kannada for a counter, then say \emph{atithigṛha} again
\end{itemize}

\begin{tcolorbox}[breakable,colback=teal!4,colframe=teal!35!black,title={Before you move on}]
What does \kn{ಅತಿಥಿಗೃಹ} mean? (\textquotedblleft{}A guest house\textquotedblright{}.) Say it once more.
\end{tcolorbox}
\section[\texorpdfstring{\kn{ಹಾಸ್ಟೆಲ್} (hāsṭel)}{ಹಾಸ್ಟೆಲ್ (hāsṭel)}]{\kn{ಹಾಸ್ಟೆಲ್} (hāsṭel) — a hostel}

\label{lesson:KA-C461-hastel}



\begin{quote}
Before the new one: say the Kannada for a counter, then the Kannada for a guest house.
\end{quote}

\subsection*{You'll want to know: \kn{ಹಾಸ್ಟೆಲ್}}
\textbf{\kn{ಹಾಸ್ಟೆಲ್}} — \emph{hāsṭel} — \textquotedblleft{}a hostel\textquotedblright{}.

\begin{quote}
\kn{ನಗರದಲ್ಲಿ} \kn{ನನ್ನ} \kn{ತಂಗಿ} \kn{ಹಾಸ್ಟೆಲ್ನಲ್ಲಿ} \kn{ಇರುತ್ತಾಳೆ}.
\end{quote}

\begin{quote}
\emph{nagaradalli nanna taṅgi hāsṭelnalli iruttāḷe.}
\end{quote}

\begin{quote}
In the city my younger sister lives in a hostel.
\end{quote}

For students, workers and budget travellers alike.

\begin{cousinweb}
\textbf{\kn{ಹಾಸ್ಟೆಲ್}} (\emph{hāsṭel}) is a loan from English \emph{hostel}.
\end{cousinweb}

\begin{grammarlens}[title={What you've built}]
One more word for getting around.
\end{grammarlens}

\subsection*{Your turn}
\begin{itemize}
\raggedright
  \item \emph{Say it:} \emph{hāsṭel}
  \item \emph{Say it:} \emph{hāsṭel}, once more
  \item \emph{Say it:} \emph{nagaradalli nanna taṅgi hāsṭelnalli iruttāḷe.} — In the city my younger sister lives in a hostel.
  \item \emph{Recall:} say the Kannada for a counter, then the Kannada for a guest house, then say \emph{hāsṭel} again
\end{itemize}

\begin{tcolorbox}[breakable,colback=teal!4,colframe=teal!35!black,title={Before you move on}]
What does \kn{ಹಾಸ್ಟೆಲ್} mean? (\textquotedblleft{}A hostel\textquotedblright{}.) Say it once more.
\end{tcolorbox}
\section[\texorpdfstring{\kn{ಠೇವಣಿ} (ṭhēvaṇi)}{ಠೇವಣಿ (ṭhēvaṇi)}]{\kn{ಠೇವಣಿ} (ṭhēvaṇi) — a security deposit; a deposit}

\label{lesson:KA-C461-thevani}



\begin{quote}
Before the new one: say the Kannada for a guest house, then the Kannada for a hostel.
\end{quote}

\subsection*{You'll want to know: \kn{ಠೇವಣಿ}}
\textbf{\kn{ಠೇವಣಿ}} — \emph{ṭhēvaṇi} — \textquotedblleft{}a security deposit; a deposit\textquotedblright{}.

\begin{quote}
\kn{ಮನೆಯ} \kn{ಬಾಡಿಗೆಗೆ} \kn{ಅವರು} \kn{ಹತ್ತು} \kn{ತಿಂಗಳ} \kn{ಠೇವಣಿ} \kn{ಕೇಳಿದರು}.
\end{quote}

\begin{quote}
\emph{maneya bāḍigege avaru hattu tiṅgaḷa ṭhēvaṇi kēḷidaru.}
\end{quote}

\begin{quote}
For renting the house they asked for ten months as a deposit.
\end{quote}

\kn{ಠೇವಣಿ} \kn{ಇಡು}, to make a deposit; \kn{ಮುಂಗಡ} is an advance.

\begin{cousinweb}
\textbf{\kn{ಠೇವಣಿ}} (\emph{ṭhēvaṇi}) is usually traced to Marathi \emph{ṭhevṇī}, \textquotedblleft{}a thing placed, a deposit\textquotedblright{}, from a verb meaning \textquotedblleft{}to put, to keep\textquotedblright{}; the path into Kannada is not settled.
\end{cousinweb}

\begin{grammarlens}[title={What you've built}]
One more word for getting around.
\end{grammarlens}

\subsection*{Your turn}
\begin{itemize}
\raggedright
  \item \emph{Say it:} \emph{ṭhēvaṇi}
  \item \emph{Say it:} \emph{ṭhēvaṇi}, once more
  \item \emph{Say it:} \emph{maneya bāḍigege avaru hattu tiṅgaḷa ṭhēvaṇi kēḷidaru.} — For renting the house they asked for ten months as a deposit.
  \item \emph{Recall:} say the Kannada for a guest house, then the Kannada for a hostel, then say \emph{ṭhēvaṇi} again
\end{itemize}

\begin{tcolorbox}[breakable,colback=teal!4,colframe=teal!35!black,title={Before you move on}]
What does \kn{ಠೇವಣಿ} mean? (\textquotedblleft{}A security deposit; a deposit\textquotedblright{}.) Say it once more.
\end{tcolorbox}
\section[\texorpdfstring{\kn{ಲಾಕರ್} (lākar)}{ಲಾಕರ್ (lākar)}]{\kn{ಲಾಕರ್} (lākar) — a locker}

\label{lesson:KA-C461-lakar}



\begin{quote}
Before the new one: say the Kannada for a hostel, then the Kannada for a security deposit.
\end{quote}

\subsection*{You'll want to know: \kn{ಲಾಕರ್}}
\textbf{\kn{ಲಾಕರ್}} — \emph{lākar} — \textquotedblleft{}a locker\textquotedblright{}.

\begin{quote}
\kn{ನಾನು} \kn{ನನ್ನ} \kn{ಪಾಸ್ಪೋರ್ಟನ್ನು} \kn{ಹೋಟೆಲ್} \kn{ಲಾಕರ್ನಲ್ಲಿ} \kn{ಇಟ್ಟೆ}.
\end{quote}

\begin{quote}
\emph{nānu nanna pāspōrṭannu hōṭel lākarnalli iṭṭe.}
\end{quote}

\begin{quote}
I put my passport in the hotel locker.
\end{quote}

Banks, stations and hotels have them.

\begin{cousinweb}
\textbf{\kn{ಲಾಕರ್}} (\emph{lākar}) is a loan from English \emph{locker}.
\end{cousinweb}

\begin{grammarlens}[title={What you've built}]
One more word for getting around.
\end{grammarlens}

\subsection*{Your turn}
\begin{itemize}
\raggedright
  \item \emph{Say it:} \emph{lākar}
  \item \emph{Say it:} \emph{lākar}, once more
  \item \emph{Say it:} \emph{nānu nanna pāspōrṭannu hōṭel lākarnalli iṭṭe.} — I put my passport in the hotel locker.
  \item \emph{Recall:} say the Kannada for a hostel, then the Kannada for a security deposit, then say \emph{lākar} again
\end{itemize}

\begin{tcolorbox}[breakable,colback=teal!4,colframe=teal!35!black,title={Before you move on}]
What does \kn{ಲಾಕರ್} mean? (\textquotedblleft{}A locker\textquotedblright{}.) Say it once more.
\end{tcolorbox}
\section[\texorpdfstring{\kn{ಮರುಪಾವತಿ} (marupāvati)}{ಮರುಪಾವತಿ (marupāvati)}]{\kn{ಮರುಪಾವತಿ} (marupāvati) — a refund; a repayment}

\label{lesson:KA-C461-marupavati}



\begin{quote}
Before the new one: say the Kannada for a security deposit, then the Kannada for a locker.
\end{quote}

\subsection*{You'll want to know: \kn{ಮರುಪಾವತಿ}}
\textbf{\kn{ಮರುಪಾವತಿ}} — \emph{marupāvati} — \textquotedblleft{}a refund; a repayment\textquotedblright{}.

\begin{quote}
\kn{ಟಿಕೆಟ್} \kn{ರದ್ದು} \kn{ಮಾಡಿದ್ದರಿಂದ} \kn{ನನಗೆ} \kn{ಮರುಪಾವತಿ} \kn{ಸಿಕ್ಕಿತು}.
\end{quote}

\begin{quote}
\emph{ṭikeṭ raddu māḍiddarinda nanage marupāvati sikkitu.}
\end{quote}

\begin{quote}
Because I cancelled the ticket I got a refund.
\end{quote}

\kn{ಮರು}-, back or again, plus \kn{ಪಾವತಿ}, a payment, as in \kn{ಪಾವತಿಸು}, to pay.

\begin{cousinweb}
\textbf{\kn{ಮರುಪಾವತಿ}} (\emph{marupāvati}) joins the native prefix \emph{maṟu}, \textquotedblleft{}again, back\textquotedblright{} (DEDR 4766), and \emph{pāvati}, \textquotedblleft{}a payment\textquotedblright{}, whose own source is not settled.
\end{cousinweb}

\begin{grammarlens}[title={What you've built}]
One more word for getting around.
\end{grammarlens}

\subsection*{Your turn}
\begin{itemize}
\raggedright
  \item \emph{Say it:} \emph{marupāvati}
  \item \emph{Say it:} \emph{marupāvati}, once more
  \item \emph{Say it:} \emph{ṭikeṭ raddu māḍiddarinda nanage marupāvati sikkitu.} — Because I cancelled the ticket I got a refund.
  \item \emph{Recall:} say the Kannada for a security deposit, then the Kannada for a locker, then say \emph{marupāvati} again
\end{itemize}

\begin{tcolorbox}[breakable,colback=teal!4,colframe=teal!35!black,title={Before you move on}]
What does \kn{ಮರುಪಾವತಿ} mean? (\textquotedblleft{}A refund; a repayment\textquotedblright{}.) Say it once more.
\end{tcolorbox}
\section[(dialogue)]{Looking back: from Planning a Trip to Guest Houses and Deposits}

\label{lesson:KA-R461-review-455-461}



\begin{quote}
Say the Kannada for tourism and a destination.
\end{quote}

\subsection*{Your turn}
Say these aloud:

\begin{itemize}
\raggedright
  \item tourism, a destination, to make a booking, a guide, worth seeing
  \item a town hall, a tower, a fountain, a landmark, to view
  \item a lane, to pass through, to explore, a spot, to look splendid
  \item a pilgrimage, a hill station, to take along, to get back, to stay overnight
  \item a mountain, a peak, to stumble, a tent, an encampment
  \item lodging, a lodge, a chamber, a receptionist, a counter
  \item a guest house, a hostel, a security deposit, a locker, a refund
\end{itemize}

\begin{tcolorbox}[breakable,colback=teal!4,colframe=teal!35!black,title={Before you move on}]
Tourism? (\textbf{\kn{ಪ್ರವಾಸೋದ್ಯಮ}}, \emph{pravāsōdyama}.) A destination? (\textbf{\kn{ಗಮ್ಯಸ್ಥಾನ}}, \emph{gamyasthāna}.) To make a booking? (\textbf{\kn{ಬುಕ್} \kn{ಮಾಡು}}, \emph{buk māḍu}.) A guide? (\textbf{\kn{ಮಾರ್ಗದರ್ಶಿ}}, \emph{mārgadarśi}.) Worth seeing? (\textbf{\kn{ಪ್ರೇಕ್ಷಣೀಯ}}, \emph{prēkṣaṇīya}.) A town hall? (\textbf{\kn{ಪುರಭವನ}}, \emph{purabhavana}.) A tower? (\textbf{\kn{ಗೋಪುರ}}, \emph{gōpura}.) A fountain? (\textbf{\kn{ಕಾರಂಜಿ}}, \emph{kārañji}.) A landmark? (\textbf{\kn{ಹೆಗ್ಗುರುತು}}, \emph{heggurutu}.) To view? (\textbf{\kn{ವೀಕ್ಷಿಸು}}, \emph{vīkṣisu}.) A lane? (\textbf{\kn{ಓಣಿ}}, \emph{ōṇi}.) To pass through? (\textbf{\kn{ಹಾದುಹೋಗು}}, \emph{hāduhōgu}.) To explore? (\textbf{\kn{ಅನ್ವೇಷಿಸು}}, \emph{anvēṣisu}.) A spot? (\textbf{\kn{ತಾಣ}}, \emph{tāṇa}.) To look splendid? (\textbf{\kn{ಕಂಗೊಳಿಸು}}, \emph{kaṅgoḷisu}.) A pilgrimage? (\textbf{\kn{ಯಾತ್ರೆ}}, \emph{yātre}.) A hill station? (\textbf{\kn{ಗಿರಿಧಾಮ}}, \emph{giridhāma}.) To take along? (\textbf{\kn{ಕರೆದೊಯ್ಯು}}, \emph{karedoyyu}.) To get back? (\textbf{\kn{ವಾಪಸಾಗು}}, \emph{vāpasāgu}.) To stay overnight? (\textbf{\kn{ತಂಗು}}, \emph{taṅgu}.) A mountain? (\textbf{\kn{ಪರ್ವತ}}, \emph{parvata}.) A peak? (\textbf{\kn{ಶಿಖರ}}, \emph{śikhara}.) To stumble? (\textbf{\kn{ಮುಗ್ಗರಿಸು}}, \emph{muggarisu}.) A tent? (\textbf{\kn{ಡೇರೆ}}, \emph{ḍēre}.) An encampment? (\textbf{\kn{ಬಿಡಾರ}}, \emph{biḍāra}.) Lodging? (\textbf{\kn{ವಸತಿ}}, \emph{vasati}.) A lodge? (\textbf{\kn{ಲಾಡ್ಜ್}}, \emph{lāḍj}.) A chamber? (\textbf{\kn{ಕೊಠಡಿ}}, \emph{koṭhaḍi}.) A receptionist? (\textbf{\kn{ಸ್ವಾಗತಕಾರ}}, \emph{svāgatakāra}.) A counter? (\textbf{\kn{ಕೌಂಟರ್}}, \emph{kauṇṭar}.) A guest house? (\textbf{\kn{ಅತಿಥಿಗೃಹ}}, \emph{atithigṛha}.) A hostel? (\textbf{\kn{ಹಾಸ್ಟೆಲ್}}, \emph{hāsṭel}.) A security deposit? (\textbf{\kn{ಠೇವಣಿ}}, \emph{ṭhēvaṇi}.) A locker? (\textbf{\kn{ಲಾಕರ್}}, \emph{lākar}.) A refund? (\textbf{\kn{ಮರುಪಾವತಿ}}, \emph{marupāvati}.)
\end{tcolorbox}
